
% \section{Discussion}
% Our work demonstrates that sustainability can be addressed at the earliest stages of the software development lifecycle by enabling energy estimation of source code. Traditional energy measurement operates at coarse granularity, limiting analysis to entire processes or long-running functions. In contrast, PowerLens enables reliable, fine-grained measurement of sub-millisecond code blocks, overcoming the temporal constraints of hardware power sensors. Coupled with our feature-based prediction model, these measurements show that energy consumption can be estimated directly from a block’s structural and syntactic properties. This capability allows developers to reason about energy usage while writing code rather than only during execution. Crucially, fine-grained measurement exposes localized energy hotspots, enabling developers to target optimization efforts where they yield the greatest benefit instead of refactoring larger code regions indiscriminately.

% By extracting energy-relevant features from the Abstract Syntax Tree (AST), our results confirm that design-time energy feedback is both feasible and practical. Integrating these predictions directly into the IDE embeds sustainability within the development workflow, providing immediate guidance similar to modern linters for style, maintainability, or security. This shifts energy optimization from a reactive practice performed after testing or deployment to a proactive design activity, supporting environmentally responsible software development from the outset. Further, the reproducibility and granularity of the measurement pipeline establish a strong empirical foundation for future research on software energy behavior.

% The approach is inherently extensible. Although trained on more than 14,000 Python code blocks, the methodology naturally scales to larger datasets, diverse library ecosystems, and additional programming languages. Because the pipeline relies on language agnostic AST structures, it can enable analyses for other languages and provides a pathway toward production-grade, multilingual energy linters.

% In prediction component, we account for the energy footprint of the predictive model itself. Energy-aware model selection requires balancing estimation accuracy with the energy overhead of generating predictions is an important trade-off, since the goal is to reduce the overall energy cost of development and execution. Our empirical results on different models enable informed further development of such models.

% To operationalize these contributions, we release PowerLens as an open-source Python library for fine-grained energy measurement and package the predictive model into a VS Code extension that provides real-time, context-specific energy estimates as developers write code. By integrating energy analysis directly into the IDE, our tools embed sustainability into day-to-day programming practice and empower developers to make informed, energy-aware decisions throughout the software development process.

\section{Discussion}
\label{sec:Discussion}
\subsection{Lessons Learned}

\textbf{Energy can now be treated as a metric of code Quality.} Several long-standing challenges identified in prior survey and vision work on energy-aware software engineering can be addressed through fine-grained, design-time energy guidance \cite{LEE2024111944}. The results for \textbf{RQ1} and \textbf{RQ3} demonstrate that energy consumption can be measured and estimated meaningfully at the level of small code blocks; this granularity aligns with how developers reason about and refactor code. In doing so, our methodology directly responds to earlier calls for improving energy observability beyond application and method-level measurements\cite{10.1145/2425248.2425252,7429297}. PowerLens reduces reliance on external hardware power meters to obtain high-resolution measurements. 

\noindent
\textbf{No single metric is a proxy for Energy.} Contrary to approaches that treat some software-quality metrics (for example, size or cyclomatic complexity) as proxies for energy\cite{10329262}, our findings for \textbf{RQ2} indicate that no single metric explains energy consumption to a large extent. Rather, energy behavior emerges from interactions among multiple code properties. The observed correlations and feature importance between block-level energy consumption and code metrics reinforce the view that energy is encoded in the compositional code features rather than isolated metrics.
The results for \textbf{RQ2} and \textbf{RQ3} suggest that energy efficient ensemble models should be used to provide actionable energy guidance. These simple machine learning approaches provide approximate, design-time estimates which can be sufficient to distinguish energy-efficient constructs from energy-intensive ones. This aligns with survey studies that emphasize the practical value of timely but approximate estimation over delayed precise measurements\cite{ZARAGOZA2025115889,LEE2024111944}. 

\subsection{Implications for Practice}
% \begin{bluepar}
From a practitioner's perspective, the key implication of this work is that energy efficiency can be considered meaningfully during code construction rather than relegated to late-stage profiling. EnCoDe's primary goal is relative hotspot ranking within a codebase, not absolute Joule-level accuracy. The classifier flags blocks predicted as \textit{High} energy tier with a lint-like warning. The developer examines the flagged block and considers alternative implementations. The regression model provides a comparative estimate between the original and refactored version to guide the choice. This workflow mirrors how developers already use static analyzers \cite{su14148596}.
% \end{bluepar}
Energy-tier classification  evaluated in \textbf{RQ3} further lowers the barrier to adoption. The system surfaces lint-like information that points to constructs likely to be energy-intensive. This mirrors successful adoption patterns for other static analyzers and suggests energy concerns can be integrated into routine development practices. Embedding energy awareness in the development lifecycle also aligns with emerging SusDevOps perspectives, where sustainability considerations are addressed alongside performance, reliability, and maintainability \cite{11023955,10.1145/3680470}. Our empirical results provide initial evidence that such integration is feasible in practice.

\subsection{Implications for Research}

For researchers, this work demonstrates how open challenges in energy-aware software engineering can be addressed via an integrated empirical and static approach \cite{LEE2024111944,ZARAGOZA2025115889,electronics14071331}. By coupling fine-grained runtime measurements with static code features, our methodology operationalizes calls to combine runtime energy data and design-time artifacts across the software lifecycle. This combined methodology enables several promising research directions: large-scale mining studies across languages and ecosystems, comparative analyzes of energy behavior across programming paradigms and systematic investigations into energy-aware refactoring.

Beyond methodological advances, our findings lend empirical support to treating energy consumption as a first-class software quality attribute. Prior work has argued conceptually for elevating energy alongside traditional qualities such as performance and maintainability; our fine-grained evidence helps close the gap by showing how energy can be measured, modeled, and reasoned about at the level of individual constructs. Finally, the availability of empirically grounded, block-level datasets opens opportunities to study energy behavior in emerging domains, such as AI-intensive systems and other compute-heavy applications, where understanding fine-grained energy signatures is increasingly critical.
